\chapter{Árboles}

\section{Definiciones Básicas}
Un \textbf{árbol} es un grafo conexo sin ciclos. Para un grafo $G$ con $n$ vértices son equivalentes:
\begin{enumerate}
  \item $G$ es un árbol.
  \item $G$ es conexo y tiene $n-1$ aristas.
  \item $G$ no tiene ciclos y tiene $n-1$ aristas.
  \item Entre cualquier par de vértices existe exactamente un camino.
\end{enumerate}

\section{Clasificación}

\subsection{Árboles binarios}
Un árbol binario es un árbol enraizado donde cada nodo tiene a lo sumo dos hijos. Propiedades:
\begin{itemize}
  \item Un árbol binario con $n$ nodos tiene $n-1$ aristas.
  \item La altura es al menos $\lceil \log_2(n+1) \rceil$.
  \item El nivel $i$ tiene a lo sumo $2^i$ nodos.
\end{itemize}

\section{Árboles Ponderados}

\subsection{Árboles generadores}
Un \textbf{árbol generador} de $G$ es un subgrafo que es árbol y contiene todos los vértices.

\subsection{Códigos de Huffman}

\begin{algorithm}[htbp]
\caption{Algoritmo de Huffman}
\begin{algorithmic}[1]
\Require Símbolos con frecuencias
\Ensure Árbol de codificación binario
\State Crear nodo hoja por símbolo con su frecuencia
\State Insertar todos los nodos en cola de prioridad
\While{cola tiene más de un elemento}
  \State Extraer los dos nodos de menor frecuencia
  \State Crear nuevo nodo con frecuencia suma
  \State Los nodos extraídos son hijos del nuevo
  \State Insertar el nuevo en la cola
\EndWhile
\State \Return raíz del árbol
\end{algorithmic}
\end{algorithm}

\section{Implementación en Haskell}

\begin{lstlisting}[caption={Árboles binarios en Haskell}]
data ArbolBinario a = Vacio | Nodo a (ArbolBinario a) (ArbolBinario a)
  deriving (Show, Eq)

preorden :: ArbolBinario a -> [a]
preorden Vacio = []
preorden (Nodo x izq der) = x : preorden izq ++ preorden der

inorden :: ArbolBinario a -> [a]
inorden Vacio = []
inorden (Nodo x izq der) = inorden izq ++ [x] ++ inorden der

postorden :: ArbolBinario a -> [a]
postorden Vacio = []
postorden (Nodo x izq der) = postorden izq ++ postorden der ++ [x]

altura :: ArbolBinario a -> Int
altura Vacio = 0
altura (Nodo _ izq der) = 1 + max (altura izq) (altura der)
\end{lstlisting}

\section{Ejercicios}
\begin{enumerate}
  \item Demostrar que un árbol con $n$ vértices tiene $n-1$ aristas.
  \item Implementar en Haskell un árbol de Huffman.
  \item ¿Cuál es la altura máxima de un árbol binario con $n$ nodos?
\end{enumerate}